% Consolidated 09-thermo-state from 7c59eacf2bbc4f2718617404cd46e24a47721d7d; original TeX fragments preserved.

% === BEGIN 09-thermo-state-001-011.tex blob b386d7a3c183d36a9817245d10121c29422330a7 ===
% Irodov, section 2.1 "Equation of the Gas State. Processes",
% original numbered problems 1--11 (printed pp.62--63; PDF pp.71--72).
% Section 2 uses separate ORIGINAL numbering; do not concatenate to 1..388.

\thermoproblem{1}{62}
一容器体积为 $V=30\,\mathrm L$，
其中装有 $0^\circ\mathrm C$ 的理想气体。
在温度保持不变的情况下放出一部分气体，
容器压强降低
$\Delta p=0.78\,\mathrm{atm}$。
气体在标准状况下的密度为
$\rho_0=1.3\,\mathrm{g/L}$。
求放出的气体质量。
\exsolution
设标准状况压强
$p_0=1.0\,\mathrm{atm}$、
温度 $T_0=273.15\,\mathrm K$。
题中容器温度同为标准温度，
因此气体的密度与压强成正比，
由理想气体状态方程，放出质量对应
\[
\Delta m=\frac{M}{RT_0}\Delta p\,V.
\]
利用标准状况
$\rho_0=Mp_0/(RT_0)$，
消去未知摩尔质量：
\[
\boxed{\Delta m=\rho_0 V\frac{\Delta p}{p_0}
=(1.3\,\mathrm{g/L})(30\,\mathrm L)(0.78)
\simeq30\,\mathrm g}.
\]
\excheck 质量与容器体积以及压强降幅成正比；
不需要单独知道气体的摩尔质量。

\thermoproblem{2}{62}
两只体积相同的容器由单向阀相连；
当两侧压强差达到或超过
$\Delta p=1.10\,\mathrm{atm}$ 时，
阀门允许气体从高压侧流向低压侧。
最初一只容器为真空，
另一只装有
$t_1=27^\circ\mathrm C$、
$p_1=1.00\,\mathrm{atm}$
的理想气体。
随后两只容器升温到
$t_2=107^\circ\mathrm C$。
求最初真空容器达到的压强。
\exsolution
两只容器总气体摩尔数不变且最终同温。
设每只容积为 $V$，
初始有气容器的气体满足
$nR T_1=p_1V$。
升温后若气体仍全部在原容器，其压强为
$p_\Sigma=p_1 T_2/T_1$。
阀门达到停止流动的临界状态时，
两侧压强之差为 $\Delta p$。
令最初真空侧压强为 $p$，
原有气体一侧压强为 $p+\Delta p$。
因为两容器体积相等、末温相同，
\[
(p+\Delta p)V+pV=nRT_2
=p_1V\frac{T_2}{T_1}.
\]
解得
\[
\boxed{p=\frac12\left(
p_1\frac{T_2}{T_1}-\Delta p\right)}.
\]
取 $T_1\simeq300\,\mathrm K$、
$T_2\simeq380\,\mathrm K$，
得到
\[
\boxed{p\simeq0.083\,\mathrm{atm}}.
\]
\excheck 原书简答记为约
$0.10\,\mathrm{atm}$；严格代入本版
$p_1=1.00\,\mathrm{atm}$、$\Delta p=1.10\,\mathrm{atm}$
给出约 $0.083\,\mathrm{atm}$。
这个差异可理解为答案过度粗略取整，
已在勘误中说明。

\thermoproblem{3}{62}
体积为 $V=20\,\mathrm L$ 的容器内
装有氢气与氦气的理想混合气，
温度 $t=20^\circ\mathrm C$、
压强 $p=2.0\,\mathrm{atm}$，
总质量 $m=5.0\,\mathrm g$。
求氢气质量与氦气质量之比。
\exsolution
令氢气、氦气质量分别为 $m_H,m_{\rm He}$，
其摩尔质量近似为
$M_H=2\,\mathrm{g/mol}$、
$M_{\rm He}=4\,\mathrm{g/mol}$。
由气体分压相加与理想气体状态方程：
\[
\frac{pV}{RT}=\frac{m_H}{M_H}
+\frac{m_{\rm He}}{M_{\rm He}},
\qquad
m_H+m_{\rm He}=m.
\]
令总摩尔数
$n=pV/(RT)$，以克为单位，
$n=m_H/2+(5-m_H)/4=(m_H+5)/4$，
所以
\[
m_H=4n-5,\quad m_{\rm He}=10-4n.
\]
取
$R=0.082057\,\mathrm{L\,atm/(mol\,K)}$、
$T=293.15\,\mathrm K$，
$n\simeq1.663\,\mathrm{mol}$，
\[
\boxed{\frac{m_H}{m_{\rm He}}
=\frac{4n-5}{10-4n}\simeq0.50}.
\]
\excheck 两组分的\emph{摩尔数}并非质量比例；
需用各自摩尔质量换算。

\thermoproblem{4}{62}
容器中混有
$m_1=7.0\,\mathrm g$ 的氮气与
$m_2=11\,\mathrm g$ 的二氧化碳，
温度 $T=290\,\mathrm K$，
总压强 $p_0=1.0\,\mathrm{atm}$。
把两气体视为理想气体，
求混合气体的质量密度 $\rho$。
\exsolution
混合气体总摩尔数等于各组分摩尔数之和：
\[
n=\frac{m_1}{M_1}+\frac{m_2}{M_2}.
\]
由理想气体方程
$p_0V=nRT$，
得到容积
$V=(m_1/M_1+m_2/M_2)RT/p_0$。
因此密度为
\[
\boxed{\rho=\frac{m_1+m_2}{V}
=\frac{p_0(m_1+m_2)}
{RT(m_1/M_1+m_2/M_2)}}.
\]
取 $M_1=28\,\mathrm{g/mol}$、
$M_2=44\,\mathrm{g/mol}$，
则两组分各为 $0.25$ mol，
总质量 $18$ g，总摩尔数 $0.50$ mol，
\[
\boxed{\rho\simeq1.5\,\mathrm{g/L}}.
\]
\excheck 混合气平均摩尔质量为
$M_{\rm av}=18/0.50=36\,\mathrm{g/mol}$，
用 $\rho=p_0M_{\rm av}/(RT)$ 可复核。

\thermoproblem{5}{63}
容器体积 $V=7.5\,\mathrm L$，
其中装有温度 $T=300\,\mathrm K$
的理想气体混合物：
氧气 $\nu_1=0.10$ mol、
氮气 $\nu_2=0.20$ mol、
二氧化碳 $\nu_3=0.30$ mol。
求：(a) 混合气总压强；
(b) 用于方程 $pV=(m/M)RT$
的平均摩尔质量 $M$。
\exsolution
气体总摩尔数
$\nu=\nu_1+\nu_2+\nu_3=0.60$ mol。
道尔顿分压定律及理想气体方程给
\[
\boxed{p=\frac{\nu RT}{V}
=\frac{0.60(0.082057)(300)}{7.5}
\simeq1.97\,\mathrm{atm}\approx2.0\,\mathrm{atm}}.
\]
混合物总质量
$m=\nu_1M_1+\nu_2M_2+\nu_3M_3$，
平均摩尔质量
\[
\boxed{M=\frac{m}{\nu}
=\frac{0.10(32)+0.20(28)+0.30(44)}
{0.60}
\simeq36.7\,\mathrm{g/mol}}.
\]
\excheck 质量加权与\emph{摩尔数加权}不同；
用组分摩尔分数计算平均摩尔质量。

\thermoproblem{6}{63}
一竖直密闭圆筒由无摩擦轻活塞分成
上下两腔，各装 1 mol 空气。
在温度 $T_0=300\,\mathrm K$ 的平衡态，
上腔体积是下腔体积的
$\eta=4.0$ 倍。
当两腔达到新的共同温度时，
体积比变为 $\eta'=3.0$，
求新的温度 $T$。
\exsolution
设圆筒总容积为 $V_\Sigma$、
活塞横截面积为 $S$、
活塞质量为 $m$。
当上、下腔体积比为 $\eta$，
\[
V_u=\frac{\eta V_\Sigma}{1+\eta},
\qquad V_d=\frac{V_\Sigma}{1+\eta}.
\]
轻活塞竖直力平衡
$(p_d-p_u)S=mg$，
理想气体状态方程
$p_d=RT/V_d$、
$p_u=RT/V_u$，
因此
\[
\frac{mg}{S}
=\frac{RT}{V_\Sigma}
\frac{\eta^2-1}{\eta}.
\]
右侧的 $mg/S,V_\Sigma$ 不变，
故
\[
\frac{T_0(\eta^2-1)}{\eta}
=\frac{T(\eta'^2-1)}{\eta'}.
\]
解出
\[
\boxed{T=T_0
\frac{\eta'(\eta^2-1)}
{\eta(\eta'^2-1)}}
=\boxed{300\cdot\frac{3\cdot15}{4\cdot8}
\simeq422\,\mathrm K}.
\]
\excheck 温度升高后两侧压强差仍由活塞重量决定，
体积比趋向更接近1；题目没有必要单独知道活塞质量。

\thermoproblem{7}{63}
体积为 $V$ 的密闭容器通过活塞式抽气泵抽气，
每次活塞行程在吸气时形成有效体积 $\Delta V$。
假设等温理想气体，
求把容器压强降低 $\eta$ 倍
所需抽气行程数 $n$。
\exsolution
一次抽气时，容器与泵室连通后
气体等温占据体积 $V+\Delta V$，
而总摩尔数暂时不变。
根据玻意耳定律，
若抽气前容器压强为 $p_j$，
抽气后重新隔离容器时
\[
p_{j+1}=p_j\frac{V}{V+\Delta V}.
\]
重复 $n$ 次：
\[
p_n=p_0\left(\frac{V}{V+\Delta V}\right)^n.
\]
要求 $p_n=p_0/\eta$，
取对数得到
\[
\boxed{n=\frac{\ln\eta}
{\ln(1+\Delta V/V)}}.
\]
\excheck 实际抽气次数须向上取整；
此式假定泵无死区容积、漏气、阀门回流或温度变化。

\thermoproblem{8}{63}
体积为 $V$ 的容器最初压强为 $p_0$，
由抽速为常数 $C$ 的真空泵等温抽气。
抽速定义为\emph{当前腔室压强下}
单位时间抽走的气体体积。
求容器压强与时间 $t$ 的关系。
\exsolution
以气体摩尔数 $n=pV/(RT)$ 描述容器状态。
单位时间抽走的摩尔数为
$\mathrm d n_{\rm out}/\mathrm dt=Cp/(RT)$。
故
\[
\frac{\mathrm d}{\mathrm dt}\left(\frac{pV}{RT}\right)
=-\frac{Cp}{RT}.
\]
因为 $T,V$ 不变，
\[
\frac{\mathrm dp}{\mathrm dt}
=-\frac CV\,p.
\]
分离变量并用初态 $p(0)=p_0$，
得
\[
\boxed{p(t)=p_0\exp(-Ct/V)}.
\]
\excheck 这里抽速 $C$ 并不是恒定\emph{抽出摩尔数}；
随着压强降低，单位时间抽出的摩尔数减少，
所以压强指数下降。

\thermoproblem{9}{63}
容积为 $V=87\,\mathrm L$ 的真空室，
抽气泵在当前压强下的抽速恒为
$C=10\,\mathrm{L/s}$。
求真空室压强降低 $\eta=1000$ 倍
所需时间。
\exsolution
根据上一题等温恒定抽速规律
$p(t)=p_0e^{-Ct/V}$。
令 $p(t)=p_0/\eta$，得
\[
e^{-Ct/V}=\eta^{-1}
\quad\Longrightarrow\quad
\boxed{t=\frac VC\ln\eta}.
\]
代入数值：
\[
\boxed{t=\frac{87\,\mathrm L}{10\,\mathrm{L/s}}
\ln1000
\simeq60\,\mathrm s\simeq1.0\,\mathrm{min}}.
\]
\excheck $V/C$ 具有时间单位；
未计入泵的极限真空度、出气和漏气。

\thermoproblem{10}{63}
原书图2.1中，一根竖直光滑管上下段横截面积不同，
上、下各有一只活塞，并用不可伸长细线相连。
两活塞之间密封 1 mol 理想气体。
上活塞面积比下活塞大
$\Delta S=10\,\mathrm{cm^2}$，
两活塞总质量 $m=5.0\,\mathrm{kg}$，
外界大气压 $p_0=1.0\,\mathrm{atm}$。
求要使两活塞整体向上平移
$\ell=5.0\,\mathrm{cm}$，气体应升温多少。
\exsolution
对两个活塞作为整体列竖直静力平衡。
活塞间气压 $p$ 对两活塞的净向上推力为
$p(S_u-S_d)=p\Delta S$，
外界大气压产生净向下力 $p_0\Delta S$，
总重力为 $mg$，故
\[
(p-p_0)\Delta S=mg
\quad\Longrightarrow\quad
p=p_0+\frac{mg}{\Delta S}.
\]
在活塞随绳共同上移 $\ell$ 时，
气体体积增加
$\Delta V=(S_u-S_d)\ell=\Delta S\,\ell$。
准静态过程中 $p$ 不变。
由一摩尔理想气体状态方程，
\[
R\Delta T=p\Delta V
=(p_0\Delta S+mg)\ell,
\]
因此
\[
\boxed{\Delta T=\frac{(p_0\Delta S+mg)\ell}{R}}.
\]
用 $\Delta S=10^{-3}\,\mathrm{m^2}$、
$p_0\simeq1.013\times10^5\,\mathrm{Pa}$、
$m=5.0\,\mathrm{kg}$、
$\ell=0.050\,\mathrm m$：
\[
\boxed{\Delta T\simeq0.90\,\mathrm K}.
\]
\excheck \textbf{题面图像复核：}
原书 PDF 纯文本抽取曾把
$\ell=5.0\,\mathrm{cm}$ 错读为 $50\,\mathrm{cm}$。
回看原版放大的题面图像，小数点清楚存在；
按 $\ell=0.050\,\mathrm m$ 计算约为
$0.90\,\mathrm K$，与本版书末简答一致。
该例\emph{不是原书答案错误}。

\thermoproblem{11}{63}
一摩尔理想气体发生下列准静态过程，
求沿该过程能够达到的最高温度，
其中 $p_0,\alpha,\beta$ 是正常数：
(a) $p=p_0-\alpha V^2$；
(b) $p=p_0e^{-\beta V}$。
\exsolution
一摩尔理想气体满足
$T=pV/R$。
(a) 温度函数
$T(V)=(p_0V-\alpha V^3)/R$，
对 $V$ 求导：
\[
\frac{\mathrm dT}{\mathrm dV}
=\frac{p_0-3\alpha V^2}{R}.
\]
内部极大值在
$V_*=\sqrt{p_0/(3\alpha)}$，
其处二阶导数
$-6\alpha V_*/R<0$。
用
$\alpha V_*^2=p_0/3$：
\[
\boxed{T_{\max}^{(a)}
=\frac{2p_0}{3R}
\sqrt{\frac{p_0}{3\alpha}}}.
\]
(b) $T(V)=p_0Ve^{-\beta V}/R$，
求导
\[
\frac{\mathrm dT}{\mathrm dV}
=\frac{p_0}{R}e^{-\beta V}(1-\beta V).
\]
唯一极大值在 $V_*=1/\beta$，
故
\[
\boxed{T_{\max}^{(b)}
=\frac{p_0}{e\beta R}}.
\]
\excheck 极值必须位于物理允许的
$p>0,V>0$ 区域：
(a) 的 $V_*$ 满足 $p=2p_0/3>0$；
(b) 所有 $V>0$ 时 $p>0$。
% === END 09-thermo-state-001-011.tex ===

% === BEGIN 09-thermo-state-012-020.tex blob a1fd6cc3a9cc6316ad8c11df7302b04049e61476 ===
% Irodov, user-supplied book, section 2.1 original problems 12--20,
% printed p.64 / PDF p.73.
\thermoproblem{12}{64}
一摩尔理想气体沿下列过程变化：
$T=T_0+\alpha V^2$，
其中 $T_0,\alpha$ 为正常数，
$V$ 为该摩尔气体的体积。
求过程中能够达到的最小压强，
并定性画出 $p$ 随 $V$ 变化的图线。
\exsolution
一摩尔理想气体满足 $pV=RT$。
代入过程方程：
\[
p(V)=\frac{R(T_0+\alpha V^2)}{V}
=R\left(\frac{T_0}{V}+\alpha V\right),
\quad V>0.
\]
对体积求导：
\[
\frac{\mathrm dp}{\mathrm dV}
=R\left(-\frac{T_0}{V^2}+\alpha\right).
\]
极小值在
$V_*=\sqrt{T_0/\alpha}$，
此点二阶导数
$2RT_0/V_*^3>0$。
于是
\[
\boxed{p_{\min}=p(V_*)
=2R\sqrt{\alpha T_0}}.
\]
图线在 $V\to0^+$ 时由 $RT_0/V$
主导而趋无穷，在 $V\to\infty$ 时
由 $R\alpha V$ 主导亦趋无穷，
在 $V_*$ 处有唯一谷底。
\excheck 过程曲线为 U 形，
不是理想气体等温过程的双曲线；
这里气体温度随体积而变化。

\thermoproblem{13}{64}
在均匀重力场中放置一个高圆柱形容器，
内装氮气。
自由落体加速度大小为 $g$，
气体沿高度 $h$ 的温度分布使得
其\emph{质量密度处处相同}。
求温度梯度 $\mathrm dT/\mathrm dh$。
\exsolution
高度向上取正，
静力学平衡给出
$\mathrm dp/\mathrm dh=-\rho g$。
由理想气体状态方程
$p=\rho RT/M$，
其中氮气摩尔质量 $M$、
$\rho$ 为质量密度。
题设 $\rho$ 在所有高度都相同，
所以
\[
\frac{\mathrm dp}{\mathrm dh}
=\frac{\rho R}{M}\frac{\mathrm dT}{\mathrm dh}.
\]
与静力方程联立：
\[
\boxed{\frac{\mathrm dT}{\mathrm dh}
=-\frac{Mg}{R}}.
\]
以氮气 $M=0.0280\,\mathrm{kg/mol}$
及 $g=9.81\,\mathrm{m/s^2}$，
得
$\boxed{\mathrm dT/\mathrm dh
\simeq-0.033\,\mathrm{K/m}=-33\,\mathrm{mK/m}}$。
\excheck 温度随高度线性下降，
与常温气体密度不随高度变化这一特定假设相容；
该梯度不是任意大气环境的实际气温直减率。

\thermoproblem{14}{64}
假设大气中的压强与质量密度满足
$p/\rho^{\,n}=\text{常数}$，
其中 $n$ 与高度无关。
利用静力学平衡和理想气体状态方程，
求相应温度梯度。
\exsolution
写多方关系为 $p=C\rho^n$，
则气体温度满足
\[
T=\frac{Mp}{R\rho}
=\frac{MC}{R}\rho^{n-1}.
\]
两边求导：
\[
\frac{\mathrm dT}{T}
=(n-1)\frac{\mathrm d\rho}{\rho},
\qquad
\frac{\mathrm dp}{p}=n\frac{\mathrm d\rho}{\rho}.
\]
由静力方程
$\mathrm dp/\mathrm dh=-\rho g$，
以及 $p/\rho=RT/M$，
得到
\[
\frac{\mathrm dT}{\mathrm dh}
=(n-1)\frac{T}{n p}\frac{\mathrm dp}{\mathrm dh}
=-(n-1)\frac{T\rho g}{np}.
\]
化简：
\[
\boxed{\frac{\mathrm dT}{\mathrm dh}
=-\frac{Mg}{nR}(n-1)}.
\]
\excheck $n=1$ 时等温且温度梯度为零；
如果 $n>1$，温度随高度升高而降低。

\thermoproblem{15}{64}
地表附近空气处于标准状况，
假设温度和平均摩尔质量随高度不变。
求地表上方 $5.0\,\mathrm{km}$ 处，
以及地面下方深 $5.0\,\mathrm{km}$ 的矿井中的气压。
\exsolution
等温静力平衡的气压分布为
\[
p(h)=p_0\exp\left(-\frac{Mgh}{RT}\right),
\]
其中高度 $h$ 在地面以上取正、矿井中取负。
取 $p_0=1.00\,\mathrm{atm}$、
$T=273.15\,\mathrm K$、
$M\simeq0.029\,\mathrm{kg/mol}$。
定义
$a=M g(5000\,\mathrm m)/(RT)\simeq0.63$。
于是
\[
\boxed{p(+5.0\,\mathrm{km})
=p_0e^{-a}\simeq0.53\,\mathrm{atm}\approx0.5\,\mathrm{atm}},
\]
\[
\boxed{p(-5.0\,\mathrm{km})
=p_0e^{a}\simeq1.9\,\mathrm{atm}\approx2\,\mathrm{atm}}.
\]
\excheck 结果来自温度在上下5 km范围内均恒定的
理想化模型；现实矿井的地温梯度与空气流动会改变数值。

\thermoproblem{16}{64}
在均匀重力场中，
假设空气温度为 $0^\circ\mathrm C$、
平均摩尔质量 $M$ 与 $g$ 均不随高度改变。
求两个高度之间的距离 $\Delta h$，
使相应空气密度分别：
(a) 相差 $e$ 倍；
(b) 相差 $\eta=1.0\%$。
\exsolution
等温大气的密度与气压同指数衰减：
\[
\rho(h)=\rho(0)e^{-h/H},
\qquad H=\frac{RT}{Mg}
\]
是密度标高。
两点间距 $\Delta h>0$ 的密度比
$\rho(h)/\rho(h+\Delta h)=e^{\Delta h/H}$。
(a) 若密度相差 $e$ 倍，
$\Delta h/H=\ln e=1$，所以
\[
\boxed{\Delta h_a=\frac{RT}{Mg}}.
\]
(b) 若相差 $1+\eta$ 倍，
\[
\boxed{\Delta h_b=\frac{RT}{Mg}\ln(1+\eta)
\simeq\eta\frac{RT}{Mg}}.
\]
取 $T=273.15\,\mathrm K$、
$M=0.029\,\mathrm{kg/mol}$，
得到
$\boxed{\Delta h_a\simeq8.0\,\mathrm{km}}$、
$\boxed{\Delta h_b\simeq0.080\,\mathrm{km}=80\,\mathrm m}$。
\excheck 当 $\eta\ll1$ 时
$\ln(1+\eta)\simeq\eta$；
不能把 $1.0\%$ 错当成倍率 $1.0$。

\thermoproblem{17}{64}
一个高为 $h$、底面积为 $S$ 的竖直圆柱容器中
装有摩尔质量为 $M$ 的等温理想气体，
温度为 $T$。
容器底面压强为 $p_0$，
重力加速度为常数 $g$。
求容器内气体总质量 $m$。
\exsolution
底端 $z=0$ 向上取高度 $z$。
等温静力平衡给出
\[
p(z)=p_0\exp(-Mg z/(RT)).
\]
理想气体质量密度为
$\rho(z)=Mp(z)/(RT)$。
因此总质量
\[
m=S\int_0^h\rho(z)\,\mathrm dz
=\frac{MS p_0}{RT}
\int_0^h\exp(-Mg z/(RT))\,\mathrm dz.
\]
积分得
\[
\boxed{m=\frac{S p_0}{g}
\left(1-e^{-Mgh/(RT)}\right)}.
\]
\excheck 对薄层高度 $h\ll RT/(Mg)$，
指数展开给
$m\simeq S p_0 M h/(RT)$，
恢复均匀近似密度 $\rho_0$ 乘体积 $Sh$。

\thermoproblem{18}{64}
把摩尔质量为 $M$ 的等温理想气体
放在均匀重力场中一个\emph{非常高}的竖直容器里，
温度为 $T$、重力加速度为 $g$。
求气体质量分布的质心高度
$h_C$（相对容器底面）。
\exsolution
等温大气密度分布为
$\rho(h)=\rho_0e^{-h/H}$，
其中 $H=RT/(Mg)$。
容器足够高，可用上限无穷大。
质心高度是质量加权平均：
\[
h_C=\frac{\int_0^\infty h\,\rho(h)S\,\mathrm dh}
{\int_0^\infty\rho(h)S\,\mathrm dh}
=\frac{\int_0^\infty h e^{-h/H}\,\mathrm dh}
{\int_0^\infty e^{-h/H}\,\mathrm dh}.
\]
由
$\int_0^\infty e^{-h/H}\mathrm dh=H$、
$\int_0^\infty h e^{-h/H}\mathrm dh=H^2$，
得到
\[
\boxed{h_C=H=\frac{RT}{Mg}}.
\]
\excheck 质心高度不受容器截面积与底部压强影响，
是这一指数密度模型的标高。

\thermoproblem{19}{64}
在均匀重力场中，理想气体摩尔质量为 $M$、
自由落体加速度为 $g$。
已知 $h=0$ 时压强 $p=p_0$，
若温度随高度分别满足
(a) $T=T_0(1-ah)$；
(b) $T=T_0(1+ah)$，
其中 $a>0$，
求相应压强分布 $p(h)$。
\exsolution
静力学平衡为
$\mathrm dp/\mathrm dh=-\rho g$，
用
$\rho=Mp/(RT(h))$ 消去密度：
\[
\frac{\mathrm dp}{p}
=-\frac{Mg}{RT(h)}\,\mathrm dh.
\]
记
$\nu=Mg/(aRT_0)>0$。
(a) 对 $T(h)=T_0(1-ah)$，
\[
\ln\frac p{p_0}
=-\frac{Mg}{RT_0}\int_0^h\frac{\mathrm dz}{1-az}
=\nu\ln(1-ah),
\]
所以
\[
\boxed{p_a(h)=p_0(1-ah)^\nu},
\qquad h<1/a.
\]
(b) 对 $T(h)=T_0(1+ah)$，
\[
\ln\frac p{p_0}
=-\frac{Mg}{RT_0}\int_0^h\frac{\mathrm dz}{1+az}
=-\nu\ln(1+ah),
\]
所以
\[
\boxed{p_b(h)=p_0(1+ah)^{-\nu}}.
\]
\excheck 两式在 $h=0$ 均给 $p_0$；
(a) 的物理定义域须保证温度 $T(h)>0$。

\thermoproblem{20}{64}
一根水平圆柱形管一端封闭、
另一端在旋转轴上与外界大气连通。
管绕穿过其开放端的竖直轴以角速度
$\omega$ 匀速旋转。
管内空气处于温度 $T$ 的等温平衡，
摩尔质量为 $M$，开放端大气压为 $p_0$。
求管内气压随距旋转轴距离 $r$ 的分布。
\exsolution
在随管转动的参考系中，
空气相对管静止。
气体在沿管轴径向满足离心静力平衡：
\[
\frac{\mathrm dp}{\mathrm dr}=\rho\omega^2r.
\]
用等温理想气体关系
$\rho=Mp/(RT)$ 代入：
\[
\frac{\mathrm dp}{p}
=\frac{M\omega^2}{RT}r\,\mathrm dr.
\]
由开放端边界 $p(0)=p_0$ 积分：
\[
\ln\frac{p(r)}{p_0}
=\frac{M\omega^2r^2}{2RT},
\]
所以
\[
\boxed{p(r)=p_0
\exp\left(\frac{M\omega^2r^2}{2RT}\right)}.
\]
\excheck 离心平衡使压强随半径\emph{增加}；
本题不是地球大气沿竖直高度衰减的模型。
% === END 09-thermo-state-012-020.tex ===

% === BEGIN 09-thermo-state-021-025.tex blob bd2f28c5753836b073d3c8f088b6d878e418882d ===
% Irodov section 2.1, original numbered problems 21--25,
% printed p.65, PDF p.74. Van der Waals equation for one mole:
% p=RT/(V-b)-a/V^2, V is molar volume (one mole).
\thermoproblem{21}{65}
在温度 $T=300\,\mathrm K$ 时，
二氧化碳的质量密度达到
$\rho=500\,\mathrm{g/L}$，
需有多大压强？
分别用理想气体方程和范德瓦耳斯方程计算。
\exsolution
设二氧化碳的摩尔质量
$M=44.0\,\mathrm{g/mol}$，
每摩尔的体积为
\[
V_m=\frac M\rho
=\frac{44.0\,\mathrm{g/mol}}{500\,\mathrm{g/L}}
=0.0880\,\mathrm{L/mol}.
\]
理想气体方程给
\[
p_{\rm id}=\frac{RT}{V_m},
\]
取 $R=0.082057\,\mathrm{L\,atm/(mol\,K)}$，
\[
\boxed{p_{\rm id}\simeq280\,\mathrm{atm}}.
\]
范德瓦耳斯方程为
$(p+a/V_m^2)(V_m-b)=RT$，
解得
\[
\boxed{p_{\rm vdW}=
\frac{RT}{V_m-b}-\frac{a}{V_m^2}}.
\]
采用该气体常用近似参数
$a\simeq3.59\,\mathrm{atm\,L^2/mol^2}$、
$b\simeq0.0427\,\mathrm{L/mol}$，
代入
\[
p_{\rm vdW}
\simeq\frac{0.082057(300)}{0.0880-0.0427}
-\frac{3.59}{(0.0880)^2}
\simeq\boxed{80\,\mathrm{atm}}.
\]
\excheck 两结果差异很大，因为该密度下
分子排除体积 $b$ 已与摩尔体积同量级，
理想气体近似失效。

\thermoproblem{22}{65}
容器体积为 $V=1.00\,\mathrm L$，
内装一摩尔氮气。
求在多高的温度下，
理想气体方程计算的压强与范德瓦耳斯方程
计算值的相对误差为 $\eta=10\%$
（相对范德瓦耳斯压强定义）；
并求此时压强。
\exsolution
取这一摩尔气体的两种压强模型：
\[
p_{\rm id}=\frac{RT}{V},\qquad
p_{\rm vdW}=\frac{RT}{V-b}-\frac a{V^2}.
\]
按本题考虑的分支，理想气体压强高于
范德瓦耳斯值，故
$(p_{\rm id}-p_{\rm vdW})/p_{\rm vdW}=\eta$，
即
$p_{\rm vdW}=p_{\rm id}/(1+\eta)$。
代入并整理：
\[
\frac{RT}{V-b}-\frac a{V^2}
=\frac{RT}{V(1+\eta)}.
\]
移项求温度：
\[
\frac a{V^2}
=RT\frac{\eta V+b}{V(1+\eta)(V-b)},
\]
\[
\boxed{
T=\frac{a(1+\eta)(V-b)}
{R V(\eta V+b)}}.
\]
此时压强
\[
\boxed{p_{\rm vdW}=\frac{RT}{V(1+\eta)}}.
\]
代入氮气的常用近似
$a\simeq1.4\,\mathrm{atm\,L^2/mol^2}$、
$b\simeq0.039\,\mathrm{L/mol}$，
得 $T$ 为约 $1.3\times10^2\,\mathrm K$，
$p_{\rm vdW}$ 为约 $10\,\mathrm{atm}$，
与书末所取参数给出的
$\sim133\,\mathrm K$、
$\sim9.9\,\mathrm{atm}$ 为相同数量级。
\excheck 题面没有提供 $a,b$ 的具体数值，
因此精确数值依赖选用的氮气材料常数表；
不可把简答中的数值视作不依赖 $a,b$ 的普适结果。

\thermoproblem{23}{65}
一摩尔某种气体装在
$V=0.250\,\mathrm L$ 的容器里。
在 $T_1=300\,\mathrm K$ 时
$p_1=90\,\mathrm{atm}$，
在 $T_2=350\,\mathrm K$ 时
$p_2=110\,\mathrm{atm}$。
求该气体的范德瓦耳斯参数 $a,b$。
\exsolution
对同一摩尔体积 $V$，
两状态均满足
\[
p_i=\frac{RT_i}{V-b}-\frac a{V^2}.
\]
相减消去 $a$：
\[
p_2-p_1=\frac{R(T_2-T_1)}{V-b},
\]
所以
\[
\boxed{b=V-\frac{R(T_2-T_1)}{p_2-p_1}}.
\]
再代回任一状态：
\[
\boxed{a=V^2\left(\frac{RT_1}{V-b}-p_1\right)}.
\]
也可代入第一式消去 $b$ 得
\[
a=V^2
\frac{p_1T_2-p_2T_1}{T_1-T_2}.
\]
取 $R=0.082057\,\mathrm{L\,atm/(mol\,K)}$，
数值为
\[
\boxed{b\simeq0.0449\,\mathrm{L/mol}},\qquad
\boxed{a\simeq1.88\,\mathrm{atm\,L^2/mol^2}}.
\]
\excheck 在本版原书简答中，
$a$ 约写成 $1.85$、
$b$ 约为 $0.042\,\mathrm{L/mol}$；
按题面印刷的两组 $p,T$ 和通用 $R$
精确代入得到上式数值。
该差别主要由取整或原版数值表造成，
公式可直接代入反查两个实验压强。

\thermoproblem{24}{65}
一摩尔范德瓦耳斯气体在温度 $T$ 下
摩尔体积为 $V$，其状态方程为
$(p+a/V^2)(V-b)=RT$。
求等温压缩系数
\[
\chi_T=-\frac1V
\left(\frac{\partial V}{\partial p}\right)_T
\]
随 $V$ 的函数。
\exsolution
首先把压强显式写为
\[
p(V,T)=\frac{RT}{V-b}-\frac a{V^2}.
\]
在温度固定时求导：
\[
\left(\frac{\partial p}{\partial V}\right)_T
=-\frac{RT}{(V-b)^2}+\frac{2a}{V^3}.
\]
故
\[
\chi_T=-\frac1V
\frac{1}{(\partial p/\partial V)_T}
=\frac1V
\frac{1}{RT/(V-b)^2-2a/V^3}.
\]
通分得到
\[
\boxed{\chi_T=
\frac{V^2(V-b)^2}
{RTV^3-2a(V-b)^2}}.
\]
\excheck 力学稳定的单相气体须满足
$(\partial p/\partial V)_T<0$，
即上式分母为正；
临界自旋odal附近压缩系数可能显著增大乃至发散。

\thermoproblem{25}{65}
接着上一题，
判断在什么温度条件下，
范德瓦耳斯气体的等温压缩系数
$\chi_{\rm vdW}$ 大于理想气体压缩系数
$\chi_{\rm id}$。
讨论摩尔体积 $V\gg b$ 时的简化结果。
\exsolution
先明确采用\emph{相同温度 $T$ 和相同摩尔体积 $V$}
比较两种状态方程。
理想气体在同一 $V,T$ 时
$p_{\rm id}=RT/V$，
等温压缩系数为
$\chi_{\rm id}=1/p_{\rm id}=V/(RT)$。
范德瓦耳斯气体在力学稳定分支的
压缩系数为
\[
\chi_{\rm vdW}
=\frac1{V[RT/(V-b)^2-2a/V^3]}.
\]
比较
$\chi_{\rm vdW}>\chi_{\rm id}$
且两个分母均为正，等价于
\[
RT>\frac{RTV^2}{(V-b)^2}-\frac{2a}{V}.
\]
移项并使用
$V^2-(V-b)^2=b(2V-b)$，
得到
\[
\frac{2a}{V}>
RT\frac{b(2V-b)}{(V-b)^2},
\]
因此
\[
\boxed{T<
\frac{2a(V-b)^2}{RbV(2V-b)}}.
\]
当 $V\gg b$ 时，
右边趋于 $a/(Rb)$，
从而
\[
\boxed{T<\frac{a}{Rb}\qquad(V\gg b)}.
\]
\excheck \textbf{与本版书末不一致：}
书末只给 $T>a/(Rb)$，方向与按\emph{同 $V,T$}
直接比较压缩系数所得的严格不等式相反。
比较不同气体的“压缩系数”必须明确
是在相同 $T,V$ 还是相同 $T,p$ 条件下；
本稿固定了前者，并将书末差异列入勘误，
不默默改写不等式方向。
% === END 09-thermo-state-021-025.tex ===
